\begin{proof}[Proof of \cref{obj:thm:finite-certificate}]
Fix the public parameters \((\beta,n,\sigma)\) satisfying the stated hypotheses. We first establish the three fixed-index guarantees simultaneously for every integer \(J\ge1\) and every \(P'\in\mathcal M_\beta(\sigma)\). For such \(J\) and \(P'\), define
\[
\begin{aligned}
A_{d,J}(P')
&=\int_0^1 K_J(x)f(P',s_dx)\mu_d(P',s_dx)\,dx,\\
B_{d,J}(P')
&=\int_0^1 K_J(x)f(P',s_dx)\,dx,\\
\eta_{d,J}(P')
&=\frac{A_{d,J}(P')}{B_{d,J}(P')},
\qquad d\in\{0,1\}.
\end{aligned}
\]
Here division by zero is assigned value zero; the positive denominator bound established below ensures that these are ordinary ratios throughout the model class. The signs \(s_1=1\) and \(s_0=-1\) are those of \cref{obj:synth_5}.

\begin{enumerate}
\item \textbf{Population identities and empirical moments.}
Define the single-observation summands, as functions on
\(\{0,1\}^2\times\mathbb R\), by
\[
\begin{aligned}
\xi^{\mathrm A}_{d,J}(D,Y,W)
&=\mathbf1\{D=d\}YQ_{J,\sigma}(s_dW),\\
\xi^{\mathrm B}_{d,J}(D,Y,W)
&=\mathbf1\{D=d\}Q_{J,\sigma}(s_dW).
\end{aligned}
\]
The inverse-heat polynomial is defined in \cref{obj:def:inverse-heat}.
Its evaluation at the proxy is square integrable: the latent score is bounded, Gaussian variables have every polynomial moment, and the indicator and binary outcome are bounded. Consequently,
\[
\xi^{\mathrm A}_{d,J},\xi^{\mathrm B}_{d,J}
\in L^2(P'_{\mathrm{obs}}).
\]

By \cref{obj:ass:gaussian,obj:ass:error-independence}, the reflected error \(s_d\varepsilon\) is standard Gaussian and independent of \((X,Y^0,Y^1)\). The first identity of \cref{obj:lem:exact-covariance} therefore gives, for every real \(x\),
\[
\int_{\mathbb R}Q_{J,\sigma}(s_dx+\sigma e)\,N(0,1)(de)
=K_J(s_dx).
\]
Integrating the independent Gaussian coordinate first, and using
\(\mathbf1\{D=d\}Y=\mathbf1\{D=d\}Y^d\), yields
\[
\begin{aligned}
\int \xi^{\mathrm B}_{d,J}\,dP'_{\mathrm{obs}}
&=\int \mathbf1\{D=d\}K_J(s_dX)\,dP',\\
\int \xi^{\mathrm A}_{d,J}\,dP'_{\mathrm{obs}}
&=\int \mathbf1\{D=d\}Y^dK_J(s_dX)\,dP'.
\end{aligned}
\]
The score marginal has density \(f(P',\cdot)\), so
\[
P'\{X=0\}=\int_{\{0\}}f(P',x)\,dx=0.
\]
Thus the arm indicator agrees almost surely with
\(\mathbf1\{s_dX\in[0,1]\}\). The defining conditional-mean identity gives the score marginal weighted by \(Y^d\) the density
\(\mu_d(P',x)f(P',x)\). Integrating the bounded arm weight against these two marginal identities gives
\[
\begin{aligned}
\int \xi^{\mathrm B}_{d,J}\,dP'_{\mathrm{obs}}
&=\int_{\{x:s_dx\in[0,1]\}}K_J(s_dx)f(P',x)\,dx
=B_{d,J}(P'),\\
\int \xi^{\mathrm A}_{d,J}\,dP'_{\mathrm{obs}}
&=\int_{\{x:s_dx\in[0,1]\}}
K_J(s_dx)\mu_d(P',x)f(P',x)\,dx
=A_{d,J}(P').
\end{aligned}
\]
For \(d=1\) the integration interval is \([0,1]\); for \(d=0\) it is \([-1,0]\), and substitution \(x\mapsto-x\) gives the displayed endpoint integrals.

For the second moments, the second identity of
\cref{obj:lem:exact-covariance}, integrated in the same order, gives
\[
\begin{aligned}
\int (\xi^{\mathrm B}_{d,J})^2\,dP'_{\mathrm{obs}}
&=\int_0^1 f(P',s_dx)
\sum_{j=0}^{m_J}\frac{\sigma^{2j}}{j!}
[K_J^{(j)}(x)]^2\,dx\\
&\le\frac34\int_0^1
\sum_{j=0}^{m_J}\frac{\sigma^{2j}}{j!}
[K_J^{(j)}(x)]^2\,dx
=V_J(\sigma),
\end{aligned}
\]
where the inequality uses \cref{obj:ass:density-bounds} and the nonnegativity of every summand. Since \(Y\in\{0,1\}\),
\[
(\xi^{\mathrm A}_{d,J})^2
=Y^2(\xi^{\mathrm B}_{d,J})^2
\le(\xi^{\mathrm B}_{d,J})^2,
\qquad
\int(\xi^{\mathrm A}_{d,J})^2\,dP'_{\mathrm{obs}}
\le V_J(\sigma).
\]

Define the centered empirical errors on the sample and seed space by
\[
\begin{aligned}
\Delta^{\mathrm A}_{d,J}(S_n,U)
&=\widehat A_{d,J}(S_n)-A_{d,J}(P'),\\
\Delta^{\mathrm B}_{d,J}(S_n,U)
&=\widehat B_{d,J}(S_n)-B_{d,J}(P').
\end{aligned}
\]
The empirical averages are those of \cref{obj:def:ratio}. For either
\(\diamond\in\{\mathrm A,\mathrm B\}\), the population identities show that
\[
\Delta^\diamond_{d,J}
=\frac1n\sum_{i=1}^n
\left[
\xi^\diamond_{d,J}(D_i,Y_i,W_i)
-\int\xi^\diamond_{d,J}\,dP'_{\mathrm{obs}}
\right].
\]
Independence makes the cross terms between distinct centered observations vanish. Expanding the square therefore gives
\[
\begin{aligned}
\mathbb E_{P',n}\Delta^\diamond_{d,J}&=0,\\
\mathbb E_{P',n}(\Delta^\diamond_{d,J})^2
&=\frac1n\int
\left(\xi^\diamond_{d,J}
-\int\xi^\diamond_{d,J}\,dP'_{\mathrm{obs}}\right)^2
dP'_{\mathrm{obs}}\\
&=\frac1n\left[
\int(\xi^\diamond_{d,J})^2\,dP'_{\mathrm{obs}}
-\left(\int\xi^\diamond_{d,J}\,dP'_{\mathrm{obs}}\right)^2
\right]\\
&\le\frac{V_J(\sigma)}n.
\end{aligned}
\]
These errors are square integrable. Integrating the independent seed contributes a factor one, since the averages depend on the sample. The same mean calculation proves
\[
\mathbb E_{P',n}\widehat B_{d,J}=B_{d,J}(P').
\]
% lean: CausalSmith.Stat.RdTruesideNoiseFrontier.marked_observed_integral_eq_endpoint
% lean: CausalSmith.Stat.RdTruesideNoiseFrontier.unmarked_observed_integral_eq_endpoint
% lean: CausalSmith.Stat.RdTruesideNoiseFrontier.marked_observed_sq_le_kernelVariance
% lean: CausalSmith.Stat.RdTruesideNoiseFrontier.unmarked_observed_sq_le_kernelVariance
% lean: CausalSmith.Stat.RdTruesideNoiseFrontier.seeded_iid_centered_moments
% lean: CausalSmith.Stat.RdTruesideNoiseFrontier.unmarkedMean_eq_single_observation

\item \textbf{Positive denominators and population bias.}
By \cref{obj:lem:positive-endpoint}, the kernel is nonnegative on \([0,1]\) and integrates to one. Hence \cref{obj:ass:density-bounds} gives
\[
\frac14
=\frac14\int_0^1K_J(x)\,dx
\le B_{d,J}(P')
\le\frac34\int_0^1K_J(x)\,dx
=\frac34.
\]
In particular, the denominator is positive and
\[
A_{d,J}(P')=\eta_{d,J}(P')B_{d,J}(P').
\]
The nonnegative weight \(K_J(x)f(P',s_dx)\), together with
\cref{obj:ass:mean-bounds}, gives
\[
\frac14B_{d,J}(P')
\le A_{d,J}(P')
\le\frac34B_{d,J}(P'),
\qquad
\eta_{d,J}(P')\in[1/4,3/4].
\]

For \(x\in[0,1]\), reflection preserves the distance from zero. Thus
\cref{obj:ass:mean-holder} gives
\[
|\mu_d(P',s_dx)-\mu_d(P',0)|\le2x^\beta.
\]
Subtracting the cutoff mean inside the weighted integral yields
\[
\eta_{d,J}(P')-\mu_d(P',0)
=
\frac{
\int_0^1K_J(x)f(P',s_dx)
[\mu_d(P',s_dx)-\mu_d(P',0)]\,dx
}{B_{d,J}(P')}.
\]
The numerator satisfies
\[
\begin{aligned}
\left|
\int_0^1K_J(x)f(P',s_dx)
[\mu_d(P',s_dx)-\mu_d(P',0)]\,dx
\right|
&\le\int_0^1K_J(x)f(P',s_dx)\,2x^\beta\,dx\\
&\le\frac32M_{J,\beta}.
\end{aligned}
\]
Since \(M_{J,\beta}\ge0\) and \(B_{d,J}(P')\ge1/4\), it follows that
\[
|\eta_{d,J}(P')-\mu_d(P',0)|\le6M_{J,\beta}.
\]
Together with the expectation identity above, this also proves the required lower bound on each unmarked population expectation.
% lean: CausalSmith.Stat.RdTruesideNoiseFrontier.endpointKernel_weighted_integral_bounds
% lean: CausalSmith.Stat.RdTruesideNoiseFrontier.arm_kernel_mean_bounds_and_bias

\item \textbf{Absolute risk.}
For every sample, define the floored denominator
\[
\widehat B^{\mathrm{floor}}_{d,J}
=\max\{\widehat B_{d,J},1/8\}.
\]
If \(\widehat B_{d,J}\ge1/8\), flooring leaves it unchanged. If
\(\widehat B_{d,J}<1/8\), then \(B_{d,J}(P')\ge1/4\) gives
\[
\left|\widehat B^{\mathrm{floor}}_{d,J}-B_{d,J}(P')\right|
=B_{d,J}(P')-\frac18
\le B_{d,J}(P')-\widehat B_{d,J}
=\left|\widehat B_{d,J}-B_{d,J}(P')\right|.
\]
Consequently, in both cases,
\[
\left|\widehat B^{\mathrm{floor}}_{d,J}-B_{d,J}(P')\right|
\le|\Delta^{\mathrm B}_{d,J}|,
\qquad
\widehat B^{\mathrm{floor}}_{d,J}\ge\frac18.
\]

Projection onto \([1/4,3/4]\) decreases distance to
\(\eta_{d,J}(P')\), which belongs to that interval. Using the population product identity,
\[
\frac{\widehat A_{d,J}}{\widehat B^{\mathrm{floor}}_{d,J}}
-\eta_{d,J}(P')
=
\frac{
\Delta^{\mathrm A}_{d,J}
+\eta_{d,J}(P')
[B_{d,J}(P')-\widehat B^{\mathrm{floor}}_{d,J}]
}{\widehat B^{\mathrm{floor}}_{d,J}}.
\]
Since \(|\eta_{d,J}(P')|\le3/4\), the projected estimate of
\cref{obj:def:ratio} satisfies
\[
|\widehat\mu_{d,J}-\eta_{d,J}(P')|
\le8|\Delta^{\mathrm A}_{d,J}|
+6|\Delta^{\mathrm B}_{d,J}|.
\]
This inequality applies to every sample, including those with an empty arm.

The triangle inequality for the two arm estimates and their population biases now gives
\[
\begin{aligned}
|\widehat\theta_J-\theta(P')|
\le{}&12M_{J,\beta}\\
&+8\bigl(|\Delta^{\mathrm A}_{1,J}|
+|\Delta^{\mathrm A}_{0,J}|\bigr)
+6\bigl(|\Delta^{\mathrm B}_{1,J}|
+|\Delta^{\mathrm B}_{0,J}|\bigr).
\end{aligned}
\]
For each of the four errors, Cauchy--Schwarz on the probability experiment gives
\[
\mathbb E_{P',n}|\Delta^\diamond_{d,J}|
\le\sqrt{\mathbb E_{P',n}(\Delta^\diamond_{d,J})^2}
\le\sqrt{\frac{V_J(\sigma)}n}.
\]
The displayed upper bound on the effect error is integrable, so integration yields
\[
\mathbb E_{P',n}|\widehat\theta_J-\theta(P')|
\le12M_{J,\beta}
+28\sqrt{\frac{V_J(\sigma)}n}.
\]
% lean: CausalSmith.Stat.RdTruesideNoiseFrontier.denominator_floor_error_le
% lean: CausalSmith.Stat.RdTruesideNoiseFrontier.clipped_ratio_error_le
% lean: CausalSmith.Stat.RdTruesideNoiseFrontier.thetaHat_absRisk_le_of_second_moments

\item \textbf{Fixed-index coverage.}
Kernel normalization implies
\[
\int_0^1K_J(x)^2\,dx>0.
\]
Indeed, if this nonnegative integral were zero, then \(K_J=0\) almost everywhere on \([0,1]\), contradicting
\(\int_0^1K_J(x)\,dx=1\). Every derivative term in \(V_J(\sigma)\) is nonnegative, and its order-zero coefficient is one, including at \(\sigma=0\). Hence
\[
V_J(\sigma)\ge\frac34\int_0^1K_J(x)^2\,dx>0.
\]
Define the positive threshold
\[
\zeta_J=\sqrt{\frac{40V_J(\sigma)}n}.
\]
The centered-moment identities imply, for every \(d\) and
\(\diamond\in\{\mathrm A,\mathrm B\}\),
\[
\operatorname{Var}_{P',n}(\Delta^\diamond_{d,J})
=\mathbb E_{P',n}(\Delta^\diamond_{d,J})^2
\le\frac{V_J(\sigma)}n.
\]
Chebyshev's inequality therefore gives
\[
\begin{aligned}
\mathbb P_{P',n}\{|\Delta^\diamond_{d,J}|>\zeta_J\}
&\le\mathbb P_{P',n}\{|\Delta^\diamond_{d,J}|\ge\zeta_J\}\\
&\le\frac{\operatorname{Var}_{P',n}(\Delta^\diamond_{d,J})}
{\zeta_J^2}
\le\frac1{40}.
\end{aligned}
\]
Define the simultaneous-control event on the sample and seed space by
\[
\Gamma_J
=\bigcap_{d\in\{0,1\}}
\left(
\{|\Delta^{\mathrm A}_{d,J}|\le\zeta_J\}
\cap
\{|\Delta^{\mathrm B}_{d,J}|\le\zeta_J\}
\right).
\]
The errors are measurable, and the union bound gives
\[
\mathbb P_{P',n}(\Gamma_J)
\ge1-\sum_{d\in\{0,1\}}
\sum_{\diamond\in\{\mathrm A,\mathrm B\}}
\mathbb P_{P',n}\{|\Delta^\diamond_{d,J}|>\zeta_J\}
\ge1-\frac4{40}=\frac9{10}.
\]
On this event the two arm errors relative to their weighted population means are each at most \(14\zeta_J\). Adding the two population biases gives
\[
|\widehat\theta_J-\theta(P')|
\le12M_{J,\beta}+28\zeta_J
=\rho_J,
\]
where \(\rho_J\) is the public radius of \cref{obj:synth_6}.

The cutoff mean bounds imply
\[
-\frac12\le\theta(P')\le\frac12.
\]
Thus, on \(\Gamma_J\),
\[
\max\{-1,\widehat\theta_J-\rho_J\}
\le\theta(P')
\le\min\{1,\widehat\theta_J+\rho_J\}.
\]
These are precisely the endpoints of \(I_J\) in
\cref{obj:synth_6}. Therefore
\[
\mathbb P_{P',n}\{\theta(P')\in I_J\}\ge\frac9{10}.
\]
Since \(J\ge1\) and \(P'\) were arbitrary, substituting \(J=L\) and \(P'=P\) proves all three fixed-index assertions.
% lean: CausalSmith.Stat.RdTruesideNoiseFrontier.kernelVariance_pos
% lean: CausalSmith.Stat.RdTruesideNoiseFrontier.centered_error_failure_le_one_fortieth
% lean: CausalSmith.Stat.RdTruesideNoiseFrontier.polyInterval_coverage_of_average_failure_bounds

\item \textbf{Honesty after public selection.}
By \cref{obj:def:public-selection}, \(L_\star\) is a deterministic function of the public parameters. Consider separately its zero and positive values.

If \(L_\star=0\), the candidate specified in \cref{obj:synth_6} is
\[
\widehat\theta_0=0,
\qquad
\rho_0=\frac12,
\qquad
I_0=[-1/2,1/2].
\]
The target bounds above give
\[
\mathbb P_{P',n}\{\theta(P')\in I_0\}=1
\qquad\text{for every }P'\in\mathcal M_\beta(\sigma).
\]
If \(L_\star\ge1\), the fixed-index coverage result gives
\[
\mathbb P_{P',n}\{\theta(P')\in I_{L_\star}\}
\ge\frac9{10}
\qquad\text{for every }P'\in\mathcal M_\beta(\sigma).
\]

For positive indices, projection of both arm estimates gives
\[
-\frac12\le\widehat\theta_J\le\frac12,
\qquad
\rho_J\ge0.
\]
Thus the centered interval and \([-1,1]\) both contain
\(\widehat\theta_J\); their intersection is a nonempty connected closed interval contained in \([-1,1]\). Its endpoints are measurable, being formed from polynomial averages, a positive denominator floor, projection, and finite maxima and minima. The zero-index interval has these properties as well. Since selection is deterministic, the selected interval has these same properties. Together with the uniform coverage bound, this proves
\[
I_{L_\star}\in\mathcal I_{\beta,n,\sigma}
\]
by \cref{obj:def:honest-class}.
% lean: CausalSmith.Stat.RdTruesideNoiseFrontier.polyInterval_properties
% lean: CausalSmith.Stat.RdTruesideNoiseFrontier.honestAttainer_of_positive_coverage

\item \textbf{Selected risk.}
If \(L_\star=0\), then for every admissible \(P'\),
\[
|\widehat\theta_0-\theta(P')|
=|\theta(P')|\le\frac12=\rho_0,
\qquad
\mathbb E_{P',n}|\widehat\theta_0-\theta(P')|\le\rho_0.
\]
If \(L_\star\ge1\), nonnegativity of \(V_{L_\star}(\sigma)/n\) and monotonicity of the square root give
\[
\begin{aligned}
\mathbb E_{P',n}|\widehat\theta_{L_\star}-\theta(P')|
&\le12M_{L_\star,\beta}
+28\sqrt{\frac{V_{L_\star}(\sigma)}n}\\
&\le12M_{L_\star,\beta}
+28\sqrt{\frac{40V_{L_\star}(\sigma)}n}
=\rho_{L_\star}.
\end{aligned}
\]
In both cases the bound is uniform over \(P'\). Since
\(E_\beta(n,\sigma)=\rho_{L_\star}\) by
\cref{obj:def:public-selection}, taking the supremum gives
\[
\sup_{P'\in\mathcal M_\beta(\sigma)}
\mathbb E_{P',n}|\widehat\theta_{L_\star}-\theta(P')|
\le E_\beta(n,\sigma).
\]
% lean: CausalSmith.Stat.RdTruesideNoiseFrontier.attainerRisk_of_positive_risk

\item \textbf{Selected expected length.}
For every positive candidate index \(J\), the endpoint formula gives samplewise
\[
\begin{aligned}
\operatorname{length}(I_J)
&=\max\!\left\{0,\,
\min\{1,\widehat\theta_J+\rho_J\}
-\max\{-1,\widehat\theta_J-\rho_J\}\right\}\\
&\le2\rho_J.
\end{aligned}
\]
Indeed, the upper endpoint is at most
\(\widehat\theta_J+\rho_J\), the lower endpoint is at least
\(\widehat\theta_J-\rho_J\), and \(2\rho_J\ge0\).
For the zero candidate,
\[
\operatorname{length}(I_0)=1=2\rho_0.
\]
Consequently the bound holds at the selected index for every sample. Integrating under the probability experiment and then taking the supremum yields
\[
\sup_{P'\in\mathcal M_\beta(\sigma)}
\mathbb E_{P',n}\operatorname{length}(I_{L_\star})
\le2\rho_{L_\star}
=2E_\beta(n,\sigma).
\]
% lean: CausalSmith.Stat.RdTruesideNoiseFrontier.attainerLength_le_certificate
\end{enumerate}
\end{proof}
